\section{Specialization}
\label{sec:spec}

The four optimizations are the four residualization cases of \Cref{sec:approach}, read off the same fixpoint. The first two need only a recovered value; the last two are the ones a scalar-injection baseline cannot reach.

\subsection{Constant propagation and finite-set specialization}
A recovered \emph{singleton} is injected as a source constant, after which the standard cascade (constant propagation, dead-code elimination, fusion, vectorization) runs unchanged. We claim no credit for the cascade; the contribution is the bridge that supplies the constant the pipeline structurally lacks (\Cref{tab:lu}). What the value unlocks is second-order and easy to undercount: once a configuration branch is known dead, the loops it separated become adjacent and fuse, so the win is the fusion, not only the removed branch. \Cref{lst:motivating} is the running case --- with \texttt{lvn\_only} and \texttt{istep} fixed, kernels \texttt{A} and \texttt{C} fuse into one column sweep --- and \Cref{sec:eval} isolates the fusion effect from the dead-code pruning that precedes it.
A recovered \emph{finite value-set} of more than one element is specialized over its members rather than folded. Each combination across several such variables is one point of a cross-product, and each point is a variant; \Cref{sec:variants} selects which variants to build.

\subsection{Index narrowing from a recovered interval}
An index array that a kernel uses only to address other arrays is bounded by the shape a configuration variable fixes. Recovering that shape as an interval $[l,u]$ lets us narrow the array's element type to the smallest integer that spans $[l,u]$, often 16 or 8 bits where the source declared 32, which halves or quarters the traffic of a gather that is frequently bandwidth-bound. This does not follow from injecting the scalar bound and re-optimizing: the storage type is fixed at declaration, and no value known later rewrites it. It also differs from static bitwidth analysis~\cite{bitwise}, which narrows only from ranges it can prove within the program; here the range is supplied by the configuration, which is what makes it available at all.

\subsection{Type-domain devirtualization}
This is the centerpiece, and the case the value-to-type reduction $\rho$ of \Cref{sec:approach} exists for. A configuration variable selects one of a finite set of concrete types, a linear solver or a tracer-advection scheme, through a discriminator, and $\rho$ turns the recovered value-set of the tag into a recovered type-set. We lower the polymorphic dispatch to an \lstinline[style=fortran]|if|-chain over that type-set, which the recovered value then prunes to the single deployed arm.
The \emph{default} production Fortran pipeline does not perform this optimization. Whole-program devirtualization is C++-vtable-specific, keyed off \texttt{!type} metadata and the \texttt{llvm.type.test} intrinsic;\footnote{\url{https://clang.llvm.org/docs/LTOVisibility.html}; \url{https://llvm.org/docs/TypeMetadata.html}.} \texttt{flang-new} emits none of that metadata and rejects \texttt{-fwhole-program-vtables}, and both \texttt{flang-new -O3} (even for a statically-known \texttt{allocate}) and \texttt{nvfortran -O4}, with or without profile-guided optimization, leave the dispatch an indirect call.\footnote{Reproduction: \texttt{experiments/devirt\_wpd/}.} The novelty is not devirtualization itself but its driver: the type-set is fixed by a configuration value behind the I/O barrier, so no compile-time analysis can resolve it without first recovering that value. $\rho$ recovers it; the devirtualization follows.

\td{VERIFY/QUESTION (skeptical review): does flang HAVE a FIR-level \texttt{fir.dispatch}$\to$\texttt{fir.call} devirtualization that is merely disabled in the default pipeline? Locate it in the flang/MLIR source and test whether enabling it devirtualizes \texttt{known.f90}. If it exists, keep the claim as stated (the DEFAULT pipeline does not, and no compile-time devirt can resolve a type fixed by a config value behind the I/O barrier); do not weaken to "no toolchain can".}
\td{For the C2 PERFORMANCE number, build an LLVM-path baseline (indirect-dispatch vs.\ devirtualized binary on the same hardware), not the DaCe "devirt vs.\ nothing-runs" comparison, which does not isolate the optimization. The reachability framing is sharpest in the dataflow backend; in an LLVM backend C2 reverts to (still novel) value-to-type dispatch elimination.}
